\section{五、进一步分析}\label{ux4e94ux8fdbux4e00ux6b65ux5206ux6790}

\subsection{（一）传导机制检验}\label{ux4e00ux4f20ux5bfcux673aux5236ux68c0ux9a8c}

借鉴江艇（2022）的建议，以渠道变量为因变量、数据要素利用及控制变量为自变量进行回归，检验信息从年报向股价传导的具体路径。三个渠道变量分别为分析师预测分歧（ForecastDisp，同一企业年度内至少3条盈利预测的标准差）、现金流波动率（CashFlowVol，经营性现金流与总资产之比的三年滚动标准差）和综合供应链集中度（SCConc，前五大客户销售占比与前五大供应商采购占比的均值）。所有渠道变量经1\%/99\%缩尾处理，模型控制企业与年份固定效应和11个控制变量，标准误在行业×年份水平聚类。

\textbf{1. 分析师预测分歧渠道}

数据要素利用涉及技术架构、业务流程和价值实现方式，投资者对其经济含义的理解存在天然门槛。当企业在年报中提供更丰富的数据利用信息时，分析师更容易围绕同一组经营事实更新预期，预测分歧因而可能下降。相较于分析师覆盖的``数量''，预测分歧衡量的是分析师群体对同一企业判断的``共识程度''，更接近信息解释是否趋于一致这一问题。

表10第(1)(2)列结果显示，DU\_kw对分析师预测分歧（ForecastDisp）的系数为-0.0065（t=-4.08），在1\%水平上显著。数据要素利用程度越高，分析师之间的盈利预测分歧越小。DU\_llm的系数不显著，表明关键词频率所承载的广度信息更有助于缩小分析师间的认知差异，而语义深度的评估尚未被分析师群体稳定吸收。考虑到该渠道仅在DU\_kw下显著，且样本量仅为22,685，该结果更适合作为补充性的信息质量证据。

\textbf{2. 经营风险渠道}

数据驱动的运营管理有助于提高需求预测精度、优化库存配置和降低生产波动，从而平滑企业经营层面的现金流。经营层面的不确定性越低，企业基本面信号越稳定，投资者对未来现金流的估值也越容易趋于收敛，信息向价格的传导因而更加顺畅。

表10第(3)(4)列结果显示，DU\_kw对现金流波动率（CashFlowVol）的系数为-0.0005（t=-2.45），DU\_llm的系数为-0.0005（t=-2.05），均在5\%水平上显著。数据要素利用降低了企业经营层面的现金流波动，双测度方向一致且均显著，经营风险渠道得到较为稳健的支持。

\textbf{3. 供应链集中度渠道}

供应链过度集中意味着企业对少数交易对手的依赖较高，一旦关键客户或供应商发生变化，经营基本面将面临更大的冲击，市场对企业价值评估的不确定性也会随之上升。数据要素利用可以提升企业对订单、库存和交易伙伴的识别与匹配能力，促进客户和供应商结构的动态优化，使供应链关系更为分散和稳定。

表10第(5)(6)列结果显示，DU\_kw对综合供应链集中度（SCConc）的系数为-0.277（t=-4.00），DU\_llm的系数为-0.460（t=-6.77），均在1\%水平上高度显著。数据要素利用显著降低了供应链集中度，双测度均高度显著，供应链渠道成立。

\textbf{表10 传导机制检验}

{\def\LTcaptype{} % do not increment counter
\begin{longtable}[]{@{}lllllll@{}}
\toprule\noalign{}
& (1) & (2) & (3) & (4) & (5) & (6) \\
\midrule\noalign{}
\endhead
\bottomrule\noalign{}
\endlastfoot
变量 & ForecastDisp & ForecastDisp & CashFlowVol & CashFlowVol & SCConc & SCConc \\
DU\_kw & -0.0065*** & & -0.0005** & & -0.277*** & \\
& (0.0016) & & (0.0002) & & (0.069) & \\
DU\_llm & & +0.0010 & & -0.0005** & & -0.460*** \\
& & (0.0017) & & (0.0002) & & (0.068) \\
控制变量 & YES & YES & YES & YES & YES & YES \\
Firm FE & YES & YES & YES & YES & YES & YES \\
Year FE & YES & YES & YES & YES & YES & YES \\
N & 22,685 & 22,685 & 43,721 & 43,721 & 42,332 & 42,332 \\
\end{longtable}
}

综合来看，现金流波动率和供应链集中度是双测度下更稳健的传导路径证据，可以视为本文机制识别的主体结果；分析师预测分歧则提供了补充性的信息质量证据。三条路径共同表明，数据要素利用不仅改变了资本市场对信息的解释方式，也与企业经营稳定性和供应链韧性的改善相联系。

\subsection{（二）广度---深度错配与概念性叙事}\label{ux4e8cux5e7fux5ea6ux6df1ux5ea6ux9519ux914dux4e0eux6982ux5ff5ux6027ux53d9ux4e8b}

前文的构念边界与增强测度检验显示，DU\_kw与更宽泛数字叙事存在交叠，但链条完整度、闭环式表述和长度标准化后的语义深度均与更快的价格发现方向一致。进一步的问题是：资本市场究竟奖励``提及得多''的概念性叙事，还是奖励``说得具体、用得更深''的实质利用披露。为此，本文构造WashGap = z(DU\_kw) - z(DU\_llm\_lenstd)，衡量关键词广度相对于语义深度的超额部分；同时构造BroadShallow虚拟变量，将当年DU\_kw位于上三分位而DU\_llm\_lenstd位于下三分位的企业识别为``提得多但说得浅''的样本。所有检验均采用与主文一致的滞后规格。

\textbf{表11 广度---深度错配与概念性叙事}

{\def\LTcaptype{} % do not increment counter
\begin{longtable}[]{@{}llll@{}}
\toprule\noalign{}
& (1) & (2) & (3) \\
\midrule\noalign{}
\endhead
\bottomrule\noalign{}
\endlastfoot
模型 & WashGap & BroadShallow & 联合回归 \\
因变量 & PriceDelay & PriceDelay & PriceDelay \\
WashGap\_lag & +0.0035** & & \\
& (0.0017) & & \\
BroadShallow\_lag & & +0.0046 & \\
& & (0.0036) & \\
DU\_kw\_lag & & & -0.0016 \\
& & & (0.0012) \\
DU\_llm\_lenstd\_lag & & & -0.0146*** \\
& & & (0.0039) \\
控制变量 & YES & YES & YES \\
Firm FE & YES & YES & YES \\
Year FE & YES & YES & YES \\
N & 37,294 & 37,294 & 37,294 \\
\end{longtable}
}

表11第(1)列显示，WashGap\_lag的系数为+0.0035（t=2.09），在5\%水平上显著为正，说明当企业在关键词广度上``说得更多''，但语义深度并未同步提高时，股价对相关信息的吸收反而更慢。也就是说，概念性叙事与实质利用之间的错配会加剧价格发现摩擦。

第(2)列中，BroadShallow\_lag的系数为+0.0046（t=1.28），方向与预期一致但统计上不显著。这意味着以分位数离散化构造的``宽而浅''虚拟变量可以提供方向一致的补充证据，但其识别强度弱于连续型错配指标WashGap。

更关键的是，第(3)列同时纳入DU\_kw\_lag与DU\_llm\_lenstd\_lag后，DU\_llm\_lenstd\_lag的系数为-0.0146（t=-3.76），在1\%水平上显著，而DU\_kw\_lag不再显著。该结果表明，在控制关键词广度之后，语义深度仍保留额外的负向解释力。结合WashGap的正向结果，以及前文DUclosedloop和DUchain\_count的负向证据，可以将这一扩展分析概括为：资本市场更倾向于奖励具有实质内容、经营嵌入和链条完整性的数据利用披露，而非单纯更频繁但较浅层的概念性提及。由于这一部分属于扩展分析，本文据此补充主文叙事，但不以其替代前文基准回归和识别结果。
